\chapter{Сети}

Контейнеры редко живут поодиночке: приложению нужна база данных,
кэш, очередь, внешние API. Docker создаёт для каждого контейнера
отдельное сетевое пространство имён и умеет соединять их в виртуальные
сети. В этой главе~--- драйверы сетей, публикация портов, разрешение
имён и диагностика типичных проблем.

\section{Как выглядит сеть контейнера}

У каждого контейнера свой сетевой интерфейс, своя таблица маршрутизации
и свой \cmd{localhost}. По умолчанию Docker подключает контейнер
к сети \cmd{bridge}, созданной на хосте. Контейнер получает внутренний
IP-адрес и может выходить в интернет через NAT, но снаружи недоступен,
пока порт не опубликован.

\begin{figure}[htbp]
\centering
\begin{tikzpicture}[
  font=\sffamily\scriptsize,
  cont/.style={draw=good!60!black, fill=good!8, rounded corners=2pt,
    minimum width=2.9cm, minimum height=0.95cm, align=center},
  bridge/.style={draw=accent!60!black, fill=accent!8, rounded corners=2pt,
    minimum width=3.6cm, minimum height=0.8cm, align=center},
  ext/.style={draw=warm!65!black, fill=warm!8, rounded corners=2pt,
    minimum width=2.4cm, minimum height=0.9cm, align=center},
  arr/.style={-Latex, thick, accent!70},
]
  \draw[draw=accent!40, fill=accent!3, rounded corners=4pt] (1.2,0.2) rectangle (12.4,5.1);
  \node[anchor=north west, text=accent!70!black] at (1.35,5.0) {хост Linux};
  \node[bridge] (br) at (6.8,3.9) {docker0 \textbullet\ 172.17.0.1/16};
  \node[cont] (web) at (4.5,1.3) {web\\172.17.0.2};
  \node[cont] (db) at (9.1,1.3) {db\\172.17.0.3};
  \draw[arr] (web.north) -- node[left, text=muted] {veth} (br.south -| web.north);
  \draw[arr] (db.north) -- node[right, text=muted] {veth} (br.south -| db.north);

  \node[ext] (user) at (0,1.3) {браузер};
  \draw[arr, warm] (user.east) -- node[above, text=warm] {localhost:8080} (web.west);
  \node[below=1pt of user, text=muted, align=center] {публикация\\8080$\rightarrow$80};
\end{tikzpicture}
\caption{Сеть bridge: контейнеры соединены виртуальным мостом,
а опубликованный порт пробрасывает трафик с хоста в контейнер}
\label{fig:bridge}
\end{figure}

Внутри контейнера \cmd{localhost}~--- это сам контейнер. Чтобы
обратиться к соседу, используется имя контейнера (в пользовательской
сети) или IP-адрес. Чтобы обратиться к сервису на хосте, нужен
адрес хоста, а не \cmd{localhost}.

\section{Драйверы сетей}

\begin{center}
\small
\begin{tabularx}{\textwidth}{@{}lXX@{}}
\toprule
\textbf{Драйвер} & \textbf{Что делает} & \textbf{Когда использовать} \\
\midrule
\cmd{bridge} & виртуальный мост на хосте, NAT & по умолчанию; локальная
разработка и большинство одиночных хостов \\
\cmd{host} & контейнер использует сеть хоста без изоляции & максимальная
производительность, но порты конфликтуют \\
\cmd{none} & только loopback, без внешней сети & изоляция, пакетная
обработка, свои сетевые настройки \\
\cmd{overlay} & сеть поверх нескольких хостов (VXLAN) & Swarm и
оркестраторы \\
\cmd{macvlan} & контейнер получает MAC-адрес в физической сети &
legacy-приложения, которым нужен L2-доступ \\
\cmd{ipvlan} & как macvlan, но на уровне L3 & высокая плотность
контейнеров \\
\bottomrule
\end{tabularx}
\end{center}

В этой книге используются \cmd{bridge}, \cmd{host} и \cmd{none};
остальные драйверы нужны в кластерах и специальных сценариях.

\section{Сеть по умолчанию и пользовательские сети}

У сети \cmd{bridge}, которую Docker создаёт автоматически, есть важный
недостаток: в ней нет автоматического разрешения имён между
контейнерами. Контейнеры могут общаться по IP, но не по имени.
Устаревший механизм \flag{--link} решал это частично и сейчас
не рекомендуется.

\begin{defbox}
\textbf{Пользовательская сеть} (\emph{user-defined bridge})~---
сеть, созданная командой \cmd{docker network create}. В ней работает
встроенный DNS: контейнеры находят друг друга по именам и сетевым
алиасам.
\end{defbox}

\begin{code}[language=bash]
$ docker network create appnet
$ docker network ls
NETWORK ID     NAME      DRIVER    SCOPE
0f1e2d3c4b5a   appnet    bridge    local
8a7b6c5d4e3f   bridge    bridge    local
...

$ docker network inspect appnet --format '{{json .IPAM.Config}}'
[{"Subnet":"172.18.0.0/16","Gateway":"172.18.0.1"}]
\end{code}

\begin{tipbox}
Создавайте отдельную пользовательскую сеть для каждого проекта.
Это даёт DNS по именам, изоляцию от чужих контейнеров и предсказуемые
имена в конфигурации.
\end{tipbox}

\section{Публикация портов}

Публикация связывает порт хоста с портом контейнера.

\begin{code}[language=bash]
$ docker run -d --name web -p 8080:80 nginx:1.27-alpine
$ docker run -d --name api -p 127.0.0.1:3000:3000 api:1.0
$ docker run -d --name dns -p 5353:53/udp coredns/coredns
$ docker run -d -P --name many nginx:1.27-alpine   # все EXPOSE на случайные порты
$ docker port web
80/tcp -> 0.0.0.0:8080
\end{code}

\begin{itemize}
  \item \flag{-p 8080:80}~--- порт 8080 на всех интерфейсах хоста.
  \item \flag{-p 127.0.0.1:8080:80}~--- только локально; снаружи
        недоступен.
  \item \flag{-p 8080:80/udp}~--- протокол UDP (по умолчанию TCP).
  \item \flag{-P}~--- опубликовать все порты из \cmd{EXPOSE}
        на случайные порты хоста.
\end{itemize}

\begin{warnbox}
\cmd{EXPOSE} в Dockerfile ничего не публикует: это документация.
Порт открывается только флагом \flag{-p} (или \cmd{ports}
в Compose). Не публикуйте базы данных на \cmd{0.0.0.0}: если порт
не нужен снаружи, привяжите его к \cmd{127.0.0.1}.
\end{warnbox}

\begin{notebox}
В Linux публикация портов реализуется правилами NAT (iptables
или nftables) и, для некоторых случаев, процессом
\cmd{docker-proxy}. Поэтому порт, занятый другим процессом,
вызывает ошибку \cmd{port is already allocated}, а изменения
в firewall хоста могут влиять на доступность контейнеров.
\end{notebox}

\section{Практика: два контейнера в одной сети}

Соединим веб-сервер и клиент, проверив DNS по имени.

\begin{code}[language=bash]
$ docker network create appnet

$ docker run -d --name web --network appnet nginx:1.27-alpine

$ docker run --rm --network appnet curlimages/curl \
    -s -o /dev/null -w '%{http_code}\n' http://web
200
\end{code}

Клиент обратился по имени \cmd{web}, и Docker разрешил его в IP
контейнера. То же самое работает с базой данных:

\begin{code}[language=bash]
$ docker run -d --name db --network appnet \
    -e POSTGRES_PASSWORD=secret postgres:17

$ docker run --rm --network appnet postgres:17 \
    psql -h db -U postgres -c 'SELECT 1;'
 ?column?
----------
        1
\end{code}

Подключение и отключение существующего контейнера:

\begin{code}[language=bash]
$ docker network connect appnet web2
$ docker network disconnect appnet web2
$ docker network connect --alias api appnet api
\end{code}

\section{Диагностика сети}

Начните с того, что видит Docker.

\begin{code}[language=bash]
$ docker network inspect appnet
$ docker inspect web --format '{{json .NetworkSettings.Networks}}'
$ docker port web
$ docker exec web cat /etc/resolv.conf
nameserver 127.0.0.11
\end{code}

Адрес \cmd{127.0.0.11}~--- встроенный DNS Docker. Если имя не
разрешается, проверьте, что оба контейнера находятся в одной
пользовательской сети, а не в сети \cmd{bridge} по умолчанию.

Дальше~--- изнутри контейнера:

\begin{code}[language=bash]
$ docker exec web getent hosts db
172.18.0.3      db
$ docker exec web ping -c 1 db
$ docker exec web curl -sS http://db:5432 || true
\end{code}

\begin{notebox}
В минимальных образах может не быть \cmd{ping}, \cmd{curl},
\cmd{getent} и даже \cmd{sh}. Для диагностики запустите временный
контейнер в той же сети:
\cmd{docker run --rm -it --network appnet nicolaka/netshoot}.
\end{notebox}

Частые причины проблем:

\begin{center}
\small
\begin{tabularx}{\textwidth}{@{}XX@{}}
\toprule
\textbf{Симптом} & \textbf{Причина и решение} \\
\midrule
\cmd{port is already allocated} & порт занят; смените порт хоста
или остановите конфликтующий процесс \\
имя не разрешается & контейнеры в разных сетях или используется
сеть \cmd{bridge} по умолчанию; создайте пользовательскую сеть \\
\cmd{connection refused} & сервис слушает \cmd{127.0.0.1} вместо
\cmd{0.0.0.0} либо ещё не готов \\
\cmd{connection timed out} &firewall или маршрутизация; проверьте
\cmd{iptables -L}, \cmd{nft list ruleset} и подсеть \\
доступ снаружи есть, с хоста нет & сервис привязан к другому
интерфейсу или порт публикуется только на \cmd{127.0.0.1} \\
\bottomrule
\end{tabularx}
\end{center}

\section{Особые режимы}

\subsection{Сеть хоста}

\cmd{--network host} отключает сетевую изоляцию: приложение слушает
порты хоста напрямую. Флаг \flag{-p} при этом игнорируется.

\begin{code}[language=bash]
$ docker run --rm --network host nginx:1.27-alpine
\end{code}

Плюс~--- нет накладных расходов NAT, минус~--- конфликты портов
и меньшая изоляция. В Docker Desktop режим host относится к виртуальной
машине, а не к вашей macOS или Windows.

\subsection{Без сети}

\cmd{--network none} оставляет только loopback. Полезно для задач
обработки данных, где сеть не нужна, и для повышения безопасности.

\subsection{Доступ к хосту из контейнера}

На Linux адрес хоста доступен через шлюз сети или
\cmd{host-gateway}:

\begin{code}[language=bash]
$ docker run --rm --add-host=host.docker.internal:host-gateway \
    curlimages/curl -s http://host.docker.internal:8080
\end{code}

В Docker Desktop имя \cmd{host.docker.internal} доступно из коробки.
Не используйте \cmd{localhost} для доступа к сервису на хосте:
внутри контейнера это он сам.

\subsection{Overlay и macvlan}

\cmd{overlay} нужен, когда контейнеры находятся на разных машинах:
сеть работает поверх VXLAN и требует кластера (Swarm или другого
оркестратора). \cmd{macvlan} выдаёт контейнеру адрес в физической
сети~--- удобно для старых приложений, которые ждут широковещательный
трафик, но требует настройки коммутатора и осторожности.

\begin{summarybox}
\begin{itemize}
  \item У контейнера своя сеть: свой \cmd{localhost}, свои порты.
  \item По умолчанию используется bridge; в пользовательской сети
        работает DNS по именам контейнеров.
  \item \cmd{EXPOSE} документирует, \flag{-p} публикует.
  \item Для диагностики смотрите \cmd{docker network inspect},
        \file{/etc/resolv.conf} и запускайте временный контейнер
        в той же сети.
  \item \cmd{host} и \cmd{none}~--- специальные режимы с другими
        компромиссами.
\end{itemize}
\end{summarybox}

\section*{Вопросы для самопроверки}
\begin{selfcheck}
  \item Почему в сети \cmd{bridge} по умолчанию не работают имена
        контейнеров?
  \item Чем \flag{-p 8080:80} отличается от \cmd{EXPOSE 80}?
  \item Как привязать опубликованный порт только к localhost?
  \item Что видит контейнер по адресу \cmd{localhost}?
  \item Как обратиться к сервису на хосте из контейнера в Linux?
  \item Почему приложение должно слушать \cmd{0.0.0.0}, а не
        \cmd{127.0.0.1}, чтобы быть доступным через \flag{-p}?
\end{selfcheck}
